\subsection{Equality in the Hopf-resolution case}

\begin{statement}{P4.27}{Proved}
Let \(S\) be an irreducible singular reduced fibre support in Part 4.2.
Suppose that the minimal resolution \(T\to N\) of its normalization has a
Hopf minimal model \(T_0\). Then
\[
 e(S)=0\quad\Longleftrightarrow\quad
 N\text{ is a smooth minimal Hopf surface}.
\]
In this case the reduced conductor and its image are disjoint unions of
smooth elliptic curves, and all component maps are finite unramified covers.
\end{statement}
\begin{proof}
Write \(\rho:T\to N\to S\), \(A=\rho^{-1}(\operatorname{Sing}S)_{\rm red}\),
and \(C=\operatorname{Sing}(S)_{\rm red}\). Let \(k\) be the number of
blowups in \(\pi:T\to T_0\). P4R.6 and P4R.8 give an effective integral pole
divisor whose reduced support \(A\) has \(r\) connected components. Each is
a smooth elliptic root with a rational tree attached. Let \(v\) be the
number of its rational components. Every such component is one of the
\(k\) blowup exceptional curves. We have
\[
 e(S)=k-v+e(C),\qquad k-v\ge0,\qquad e(C)\ge0.
\]
Thus equality forces \(k=v\) and \(e(C)=0\).

Let \(c,s,g,b\) denote the numbers of irreducible components of \(C\),
connected components of \(C\), elliptic normalized components of \(C\), and
the sum of its point-identification graph cycle ranks. P4R.8 gives
\[
 e(C)=c+s-2g-b,\qquad
 b+r-s\le1,\qquad c\ge g,\quad r\ge g.
\]
Consequently \(e(C)=0\) implies \(c+r-2g\le1\).

First suppose \(c=g\). Every normalized component of \(C\) is elliptic.
Every rational component of \(A\) is contracted by \(T\to N\), since a
finite nonconstant map from \(\mathbb P^1\) to an elliptic normalization
does not exist. A root cannot be contracted. Indeed all rational curves
attached to it are already contracted, so its whole rooted connected
component would be a connected exceptional block disjoint from the strict
conductor. P4R.6 excludes precisely such a block. Hence all roots are
noncontracted.

The equality \(k=v\) now says that every blowup exceptional curve of
\(\pi:T\to T_0\) belongs to \(A\) and is contracted by \(T\to N\).
Conversely every exceptional curve of \(T\to N\) belongs to \(A\), by
P4R.6, and cannot be a root. The two proper birational maps therefore have
exactly the same exceptional curves. Their fibres are identical: each
nontrivial fibre is a connected component of the exceptional locus, and
their fibres are connected by Stein factorization and normality of the
targets.

Here is the holomorphic comparison, rather than only a topological
identification. Take the image of
\[
 (\pi,\rho_N):T\longrightarrow T_0\times N
\]
and then its normalization \(G\). The two projections \(G\to T_0,N\)
are proper birational maps with zero-dimensional fibres, because the
original two maps have identical fibres. They are finite. A finite
birational morphism to a normal complex space is an isomorphism, as follows
on local rings from integral closedness. Thus \(N\simeq T_0\)
biholomorphically. It is a smooth minimal Hopf surface.

The only other possibility is \(c=g+1,r=g\). There is one rational
normalized component of \(C\), and all \(g=r\) roots are noncontracted and
map to distinct elliptic normalized components of \(C\). Each such map is
an isogeny. Also
\[
 b+r-s=1.
\]
We exclude this possibility by a quotient of the exact fundamental-group
pushout; the quotient is constructed explicitly below.

Represent \(C\) as its disjoint normalized smooth curves with a finite
graph of point identifications. This is a graph of spaces with elliptic
vertex groups \(\mathbb Z^2\), trivial rational vertex groups, and trivial
edge groups. Its graph has \(s\) connected components and total cycle rank
\(b\). Add a vertex for \(T\), and add one root edge for each connected
component of \(A\). The root edge group is \(\mathbb Z^2\), since its
rational tree is simply connected. Its map to the \(C\) side factors
through the elliptic normalization receiving that root; its map to the
\(T\) side is the actual inclusion. This graph-of-spaces presentation is
the ordinary van Kampen presentation of the homotopy pushout
\(T\cup_A C=S\). It has total graph cycle rank \(b+r-s=1\).

Let \(\Gamma=\pi_1(T)=\pi_1(T_0)\). It is virtually infinite cyclic, and
\(b_1(T)=1\). The primitive map from its abelianization to its free part
therefore gives a surjection
\[
 h:\Gamma\longrightarrow\mathbb Z
\]
with finite kernel. Indeed a cyclic subgroup of finite index maps
nontrivially to \(\mathbb Z\), so the kernel meets that subgroup trivially
and is finite.

For root \(H_i\), write
\[
 a_i:\pi_1(H_i)=\mathbb Z^2\longrightarrow\Gamma,\qquad
 q_i:\pi_1(H_i)=\mathbb Z^2\longrightarrow
       \pi_1(\widetilde C_i)=\mathbb Z^2.
\]
The second map is injective of finite index because it is induced by an
isogeny. The composite \(h a_i\) is nonzero. To see this, a finite image
of \(a_i\) would become trivial after a finite cover of \(H_i\); that
compact elliptic cover would lift holomorphically to the universal Hopf
cover \(\mathbb C^2\setminus\{0\}\), where its coordinate functions are
constant. Thus \(a_i\) has infinite image; since \(\ker h\) is finite,
\(h a_i\ne0\).

Put \(K_i=\ker(h a_i)\), a primitive rank-one subgroup of
\(\mathbb Z^2\). Quotient the \(i\)-th elliptic vertex group by the
saturation of \(q_i(K_i)\). The resulting vertex group is \(\mathbb Z\).
Quotient the root edge group by \(K_i\); it too becomes \(\mathbb Z\).
Its two induced maps into the \(T\) vertex \(\mathbb Z\) and the elliptic
vertex \(\mathbb Z\) are nonzero multiplications and hence injective.
The latter assertion follows because an element outside \(K_i\) cannot
map under the injective rational map \(q_i\otimes\mathbb Q\) into the
rational span of \(q_i(K_i)\). This verifies compatibility of the chosen
quotients with both edge homomorphisms.

We have constructed a surjective homomorphism from \(\pi_1(S)\) onto the
fundamental group \(G\) of a graph of groups. Its nontrivial vertex groups
are infinite cyclic, its root edges are infinite cyclic with injective
maps, and all point-identification edges are trivial. The root edges
form a star: the roots map to distinct elliptic normalizations. They
therefore contain no graph cycle. The unique graph cycle of the complete
graph contains a point-identification edge, whose edge group is trivial.

Delete one such edge on the cycle. The remaining graph is a connected
tree, and the normal form for iterated amalgamated products along
injective edge maps embeds its \(T\) vertex group \(\mathbb Z\) into its
fundamental group \(G_0\). Restoring the removed edge is an HNN extension
over the trivial group, hence
\[
 G\simeq G_0*\mathbb Z.
\]
Choosing a nontrivial element of the embedded \(\mathbb Z\subset G_0\)
and the new stable letter gives an embedded free group \(F_2\), by reduced
free-product normal form. But every quotient of a virtually abelian group
is virtually abelian, and a virtually abelian group has no \(F_2\)
subgroup. This contradicts P4R.1. The case \(c=g+1,r=g\) is impossible.

We have proved that equality makes \(N\) smooth minimal Hopf. Conversely
P4.24 applies to such an actual normalization and proves the unramified
elliptic-conductor description and \(e(S)=0\). This completes both
directions.
\end{proof}
